% Notes of the Montreal seminar, 1970 (cote 129, typed part): archive cote 129, site page 47 (PDF page 48).
% typescript. Draft compared: none.
% AI-transcribed by Claude Opus 5.5 (claude-opus-5-5, adaptive thinking, effort high), Message Batches API, batch mtl129_full (msgbatch_01AGnDKHwy6ATCM7ehUa6263), 2026-11-13; not yet checked by eye.
\scanpage{129/47}
\mspage{\hand{II 19}}

\margin{$\times$}
\[
\underline{\mathrm{Mod}}_{\hand{$\mathcal{D}$};W\,\hand{$\ell.f.$}} \xrightarrow{\ \approx\ }
\left(\begin{array}{l}
\text{sous-catégorie pleine}\\
\text{de } \underline{\mathrm{Mod}}_{\hand{$\mathcal{D}$};W\,\hand{$\ell.f.$}} \text{ formée}\\
\text{des objets où } V \text{ est}\\
\text{nilpotent}
\end{array}\right)
\mathrm{x}
\left(\begin{array}{l}
\text{sous catégorie pleine}\\
\text{de } \underline{\mathrm{Mod}}_{\hand{$\mathcal{D}$};W\,\hand{$\ell.f.$}} \text{ formée}\\
\text{des objets où } V \text{ est}\\
\text{bijectif}
\end{array}\right) .
\]

Cela résulte du lemme suivant :

\uline{Lemme}. Tout \hand{$\mathcal{D}$}-module $M$ de longueur finie sur $W$ se décompose d'une
manière unique (et fonctorielle) en une somme directe
\[
M = M' \;\mathrm{x}\; M''
\]
telle que $V|M'$ est nilpotent et $V|M''$ est bijectif.

\uline{Preuve}. Soit $N$ un entier suffisamment grand pour que
\[
\bigcup_{n\underline{>}0} \mathrm{Ker}\, V^{n} = \mathrm{Ker}\, V^{N} \quad \text{et} \quad \bigcap_{n\underline{>}0} \mathrm{Im}\, V^{n} = \mathrm{Im}\, V^{N} .
\]
On a évidemment $V|\mathrm{Ker}\, V^{N}$ nilpotent et $V|\mathrm{Im}\, V^{N}$ bijectif et de plus
$M = \mathrm{Ker}\, V^{N} \oplus \mathrm{Im}\, V^{N}$ . Il est clair que cette décomposition est unique et
fonctorielle.

5. \uline{Corollaires et compléments}

\uline{Corollaire 5.1}

Le foncteur $D^*$ commute aux changements de base définis par
toute extension parfaite $K$ de $k$ : Pour un objet $G$ de $p$-Gr-fin$(k)$ ,
on a un isomorphisme fonctoriel
\[
D^*(G_{K}) \overset{\sim}{\leftarrow} W(K) \otimes_{W(k)} D^*(G) .
\]

\mspage{\hand{47}}
